\section{Assumptions and limitations}
\label{sec:assumptions}

\subsection{Assumptions}
Table~\ref{tab:assumptions} lists what the relations take for granted, and what
each assumption costs in accuracy.

\begin{table}[!ht]
  \centering
  \small
  \renewcommand{\arraystretch}{1.25}
  \begin{tabularx}{\linewidth}{@{}l X@{}}
    \toprule
    \multicolumn{1}{c}{\textbf{Assumption}} & \textbf{Consequence and justification} \\
    \midrule
    Purely resistive DC line
      & No inductance or capacitance appears in the relations. Correct for a
        steady-state DC flow; line charging and inductance matter only for
        transients and fault studies. \\
    Uniform conductor temperature
      & The whole route sits at the single design temperature $T$. In service
        the temperature varies with ambient, wind and loading along the line,
        so the reported loss is that of the stated design condition rather
        than of any particular hour. \\
    Linear resistance-temperature law
      & Equation~\eqref{eq:ktemp} is the standard inferred-zero-temperature
        form, accurate to well within a percent from ambient to the maximum
        continuous rating of an ACSR conductor. \\
    Aluminium carries the current
      & The steel core of an ACSR conductor is neglected in the resistance,
        which is standard practice; $T_{0} = 228\,^{\circ}$C is the aluminium
        constant. \\
    Station losses proportional to throughput
      & No no-load term, so light-load efficiency is overstated while
        full-load efficiency is representative.\\
    Balanced bipole carries no DMR current
      & The pole currents are taken as exactly equal. In service a small
        unbalance flows in the DMR; its loss contribution is second order,
        going with the square of a few percent of $I_{d}$. \\
    \bottomrule
  \end{tabularx}
  \caption{Assumptions underlying the load flow relations.}
  \label{tab:assumptions}
\end{table}

\subsection{Outside the scope}
The following are deliberately not treated here, each belonging to a dedicated
study:

\begin{itemize}\setlength{\itemsep}{0.3em}
  \item \textbf{Reactive power and AC-side voltage.} Only active power crosses
        the converter boundary in these relations. Converter reactive
        capability and AC network voltage support belong to the AC side of the
        design and, ultimately, to interconnection studies.
  \item \textbf{Thermal rating of the conductor.} The flow yields the current
        and assumes a conductor temperature; whether that current actually
        produces that temperature at the ambient and wind conditions of the
        route is an ampacity calculation, and the two should be reconciled
        once the route conditions are fixed.
  \item \textbf{Dynamics.} Control response, ramp rates, pole-block
        transients and the transition into the one-pole-out state are all
        outside a steady-state flow.
\end{itemize}

\subsection{Open items}
\begin{enumerate}[label=(\arabic*)]
  \item \textbf{Conductor temperature basis.} 75\,$^{\circ}$C is the maximum
        continuous temperature and gives a worst-case loss and voltage drop.
        Whether losses should also be reported at an average operating
        temperature is a reporting decision still open --- the two differ by
        of the order of 10\,\% in loss.
  \item \textbf{Converter loss figure.} 1.2\,\% per station is a generic
        VSC-class value. To be replaced with vendor loss curves, at which
        point the two-term loss model of Table~\ref{tab:assumptions} becomes
        worthwhile.
\end{enumerate}
